\section{Future Work}\label{sec:future-work}
The developed results constitute a reformulation of established physics in
space-time without gravitational curvature. Extensions not yet constructed
are separated from points requiring
greater mathematical precision. A proposed extension is marked
\textbf{TODO}. A gap requiring resolution before physical claims can be
supported is marked \textbf{FIXME}. No premise for the established results is
taken from this section.

\subsection{TODO: Local Frames and Gravitation}
One fixed sigma basis is used throughout space and time in the present calculations.
A gravitational extension would instead need a position-dependent local
basis and a rule for comparing algebra elements defined at neighboring
positions.  That comparison rule would have to distinguish an ordinary
change of description from a physical gravitational field.  In particular,
a basis written everywhere as a single transform of a fixed basis is locally
only a change of description. A failure to return an algebra element
unchanged after comparison around a sufficiently small closed path is called
curvature \cite{doran2003}; it is not produced by that basis
change alone.

A TODO question is therefore whether one comparison rule can be defined and
used consistently in every equation. Such a rule is conventionally called a
connection \cite{doran2003}. Curvature would then have to be calculated by
the closed-path comparison just defined, rather than by a change of basis
alone. If two successive comparisons depend on their
order, they are said not to commute. The interval, freely falling motion,
Maxwell equation, Dirac system, and energy--momentum source rule would all
have to be rebuilt with that same comparison rule. Agreement with the
constant-basis formulas would be a required check.

It remains to be determined whether translational motion can be combined
with rotation-rate and boost-rate cross-position terms and consistently
interpreted as a velocity. Multiplication by mass and construction of a
transformed interval would follow only if that interpretation were proved.
Those physical identifications are not justified by the moving-basis generator derived in
Sec.\ \ref{sec:transform-algebra} alone. In particular, surrounding each
basis element by a transform and its inverse is not the same operation as the
two-sided event transformation used for a Lorentz boost, and a complex
velocity-like expression is generally produced by the proposed boost-rate
term. It must therefore be determined whether a real event law and an
observer's measured time and distance can be defined, whether an ordinary
derivative can be corrected for the changing local basis so that the
corrected derivative transforms with the field it acts on, and whether real
measured velocity, energy, momentum, and interval can then be derived. The
corrected derivative is conventionally called a covariant derivative
\cite{doran2003}. Recovery of the inertial formulas when the
basis-change rate is constant or zero would be required.

Space-time curvature is also not yet supplied by a plane-curve analogy.
A numerator obtained by differentiating a coordinate quotient is only one
piece of a curve formula. The remaining TODO questions are whether the curve
parameter has a fixed physical meaning, whether the vector pointing along
the curve, conventionally called its tangent, has a fixed length, how
directions are compared at neighboring positions, and whether the proposed
result works for motion in more than one direction. Curvature
cannot be inferred from that numerator alone.

\subsection{TODO: Additional Field-Value Transformations}
One phase freedom and one potential paravector are used in the electromagnetic
coupling. Here an internal field value means a label carried by the field
that is changed without changing the space-time event at which the field is
evaluated. Other established field theories use several coupled internal
values and transformations for which applying the first and then the second
can differ from applying them in the opposite order \cite{doran2003}. The
TODO question is whether such transformations can be added while keeping the
space-time product and the internal action separate.

The allowed internal field values, their transformation rule, the
corresponding potentials, and the quantities preserved by the resulting
equations would have to be defined before any field equation is written.
The products applied to space-time paravectors and those applied only to the
internal field values would also have to be shown explicitly. It must be
determined whether the existing algebra can express both actions without
ambiguity. Merely replacing the electromagnetic potential by a larger array
would not answer that question.

\subsection{TODO: Many-Particle and Particle-Number-Changing Fields}
Classical fields and single-particle wave functions are described by the wave
equations in this manuscript. The term ``first quantized'' refers to this use
of an ordinary wave function while the number of particles remains fixed. In
particular, a negative-frequency solution is not by itself an antiparticle
\cite{greiner2000}. A future theory would require field operators whose
action can change the number of particles, operations that add or remove one
particle, a transformation that reverses charge, and a no-particle reference
state called the vacuum. The TODO question is whether all of these objects
can be represented without obscuring the present algebra, and whether the
particle--antiparticle interpretation can then be derived rather than
assumed \cite{peskin1995}.

Further TODO questions concern how the state is changed when two identical
particles are exchanged, how a state is represented when it cannot be
written as one independent state for each particle, and how the remaining
state is described after some particles are left unmeasured. The condition
that the state cannot be separated is conventionally called entanglement,
and the last partial
description is conventionally called a reduced state. The rule for extracting
quantities measured near one position would have to be specified in the same
development \cite{peskin1995}.
The fixed projected column space used for one Dirac field is a useful
starting point, but these many-particle constructions are not supplied by it
alone.

The indefinite Klein--Gordon pairing and the absence of a positive local
one-particle position density are also not repaired by a change of algebraic
packaging. In a theory with particle-number-changing field operators, the
origin of a positive squared state length, an operation that measures charge,
and the distinction between particles and antiparticles must be shown. The
one-particle formulas must then be recovered when particle-number changes can
be neglected.

\subsection{TODO: Electromagnetic Response of Materials}
Empty-space fields and prescribed sources are treated in the electromagnetic section.
A material requires an additional rule that converts the electromagnetic
field into the response of that material \cite{jackson1998}. Such a rule is
conventionally called a constitutive map. Its outputs can include electric
dipole moment per unit volume, called polarization; magnetic dipole moment
per unit volume, called magnetization; and electric current. An electric
dipole moment records the amount and direction of separated positive and
negative charge. A magnetic dipole moment records the strength and direction
of a small magnet or current loop. Direction
dependence is called anisotropy, frequency dependence is called dispersion,
current produced by an electric field is described by conductivity, and
conversion of field energy into heat or other unavailable forms is called
dissipation.

The first TODO question is whether simple media whose response respects
addition and multiplication by an ordinary number can be written clearly.
Such a response rule is called linear. Direction-dependent and
time-dependent responses can then be tested. The distinction between a field
paravector and an operation acting on that field must be preserved, and the
energy balance must be derived rather than assumed. It must then be
determined which material laws fit into one paravector product and which
require a separately defined map.

\subsection{FIXME: Operator Foundations}
Differential, inverse, and square-root operators are used in several
factorizations in the wave-mechanics sections. Their algebraic
forms are correct on the stated test fields, for which all displayed
derivatives exist, but the allowed set of fields must be specified
completely. That allowed set is the operator's domain. Conditions imposed at
the edge of the spatial region or at infinite distance are its boundary
conditions \cite{greiner2000,sakurai2017}. The conditions under which
integration by parts---moving a derivative from one factor in an integral to
the other factor---adds no boundary term, an inverse exists, and an energy
operator has real measured values are determined by those choices.

A consistent definition of the square-root frequency operator is also
required for every plane-wave term. The nonnegative square root must be
selected explicitly, and the spatially constant term with
\(\vct{k}=\vct{0}\) must be treated separately when the inverse is used.
Enough derivatives and integrals must exist for both written operation orders
in every use of the commutator function. These conditions should eventually
be stated once in a common operator-domain section and then cited wherever
they are used in a factorization or low-energy expansion.

Likewise, the determinant of a block differential operator minus a proposed
measured value is not defined merely because the paravector determinant is
defined. An ordinary matrix determinant whose zeros select possible values
may be used only after the equation has been reduced to a finite matrix, for
example one plane wave at a time, a finite set of basis functions, or a
finite numerical grid. Merely restricting the spatial region does not make
the operator finite-dimensional. At the operator level, only the ordered
factorizations and explicitly proved
polynomial identities are presently justified. It remains to be determined
under what stated conditions an infinite-dimensional determinant can be
defined, if one is needed at all.

For scalar relativistic waves, the positive-energy solution space for each
permitted background that does not change with time must be constructed, and
the measurement operators that preserve it must be determined. General
time-dependent backgrounds can mix the positive- and negative-frequency
parts, and a position probability cannot be obtained merely by declaring the
local Klein--Gordon density to be a position-probability density. For Dirac
and Pauli wave-function
spaces, the shared operator domains must be specified. It must then be proved
that moving an operator from one side of the bracket to the other gives the
displayed \(H\) operation, that the energy-selecting operators preserve their
selected states, and that the calculated measurement averages are real.

\subsection{FIXME: Bounds on Measurement Spread}
A spin-measurement uncertainty bound cannot be inferred directly from the
axis commutator function and a determinant inequality. For a measured
quantity, its centered operator is obtained by subtracting its average from
the operator. Its variance is the average square of that difference, and its
uncertainty is the nonnegative square root of the variance
\cite{robertson1929,sakurai2017}. The normalized projected-wave bracket in
Eq.\ \eqref{eq:adjoint-paravector-braket} must be used to define those
averages. A shared set of fields on which both operator orders and their
\(H\) operations exist must also be required. The standard lower bound built
from the commutator function, and the stronger bound that also retains the
anticommutator function, should then be proved in the manuscript's notation.
Spin-axis measurement operators should be checked as a worked example rather
than treated as the proof.

\subsection{TODO: Higher-Order Pauli Reduction}
The leading kinetic and magnetic spin terms are proved by the present Dirac
reduction. The next TODO question is whether the next-order terms can be
derived with the same sign and charge conventions. These include the coupling
between spin and motion in an electric field, called the spin--orbit term,
and the local correction produced by spatial variation of the electric
field, often called the contact or Darwin term
\cite{greiner2000,sakurai2017}. Corrections required for time-dependent
potentials must also be derived. The density and current must be expanded to
the same order as the energy operator; retaining higher-order energy terms
while leaving the measurement rule at leading order would be inconsistent.
These terms are established consequences of the Dirac equation, but their
derivation in this notation remains future work.

\subsection{TODO: Classify Companion Components}
Paired real and imaginary components and paired scalar and vector components
are repeatedly produced by the algebra. Some of these are familiar physical
laws: the magnetic-divergence and induction equations, for example, are
encoded by the imaginary parts of the Maxwell equation. Other components, such as the
imaginary part of the force product, are valid algebraic companions without
a separate physical interpretation established here.

Two cases are already identified with established physics: a projected
block's opposite-order product is the standard density matrix for one fully
specified two-entry state \cite{sakurai2017}. Its rank is one, meaning that
every column is a multiple of the one selected state column. The remaining
\(\pv{F}\) and \(\conj{\pv{F}}\) terms in the two ordered products are the
standard two Dirac spin couplings \cite{greiner2000}. For the remaining
nonscalar channels of unrestricted or differently ordered products, a
separate derivation as a current, spin-axis quantity, or measurement result
is still required before physical meaning is assigned. For a fixed-column
state and a left-acting operator that preserves that column, the product
\(\herm{\pv{\psi}}\pv{\mathcal{O}}\pv{\psi}\) is already fixed by its one
complex row--column matrix element and contains no additional independent
channels.

The following exact identity may also be studied:
\begin{equation*}
\Box(\pv{F}^{\,2})
=\conj{\pv{\partial}}\pv{\partial}(\pv{F}^{\,2}).
\end{equation*}
Because \(\pv{F}^{\,2}\) is a complex scalar, this equality is an exact use
of the defined wave operator, but it is not an independent field equation.
Any useful source expansion must be derived from Maxwell's equation with the
ordered product rule of Eq.\ \eqref{eq:paravector-leibniz}. A TODO question
is whether the resulting scalar identity describes a balance unchanged by
the stated transformations or merely repeats information already contained
in Maxwell's equation. It must not be presented as an independent law of
field evolution without such a derivation.

Each companion formula should be placed by a systematic classification into
one of three categories: an alternate form of an established law, redundant
information implied by another displayed equation, or an algebraic channel
whose physical meaning is presently unassigned. The operations that return
the original input when applied twice, the order functions,
\(\pv{\Pi}(\pm)\) selection functions, and field products should all be
included in that classification. No unassigned
component should be advertised as a new
interaction merely because it is made available by the algebra.
Maps that preserve the defined addition and multiplication, together with
the selection, factorization, exponential, and logarithmic constructions,
should be separated into those proved directly from the algebra and those
that depend on the chosen matrix realization.

\subsection{TODO: Complex Space-Time and Imaginary Coordinates}
A complex paravector \(\pv{Z}\in\Pspace(\C)\) is already permitted by the
generic algebra: \(\pv{Z}=\pv{X}+i\pv{Y}\), with
\(\pv{X},\pv{Y}\in\Pspace(\R)\). Treating such an element as a physical event, rather than as an
algebraic field value, would be a new hypothesis. The following self-contained
consistency identity is obtained from the existing determinant and adjugate
definitions:
\begin{equation*}
\detf{\pv{Z}}
=\detf{\pv{X}}-\detf{\pv{Y}}
+i\Bigl(
  \pv{X}\adjf{\pv{Y}}
 +\pv{Y}\adjf{\pv{X}}
\Bigr).
\end{equation*}
The quantity in parentheses is a scalar, but the determinant of
\(\pv{Z}\) is generally complex. No interval interpretation is fixed by this
identity. In a proposed complex-event theory, it would have to be stated which,
if any, real quantity is measured as an interval and why that choice is
preserved by the allowed transformations.

The scalar part of \(\pv{Y}\) could be described as an imaginary-time
coordinate, and its vector part could be described as imaginary space.
Those descriptions alone carry no physical content. A metric is a stated
rule that assigns an interval to a coordinate difference
\cite{minkowski1908}. The TODO question is whether a real interval rule can
be selected for a complex paravector and preserved by the allowed
transformations. Taking the real part, absolute magnitude, or squared
absolute magnitude of a complex determinant would give different interval
rules, so none can be selected by terminology alone. The ordinary real
interval would have to be recovered when the imaginary part vanishes.

Only after that choice is justified could the time measured along an
observer's path, the possible time order of events connected by a signal,
velocity, momentum, or force be defined for a complex event. For real
space-time, the zero-interval directions form the light cone and separate
event differences that can be joined by a signal no faster than light from
those that cannot \cite{einstein1905,minkowski1908}. A TODO question is
whether that same measured boundary is recovered to experimental accuracy.
Complex mass, unobserved paths, or new force channels should not be introduced
before these more basic questions are resolved.

Faster-than-light propagation is not implied by complex coordinates. A
propagation speed can be assigned only after field equations, the surfaces
along which disturbances propagate, observers, and measured time and distance
have been defined. If
the measured projection remains the real paravector \(\pv{X}\), the ordinary
light cone may remain unchanged. If \(\pv{Y}\) changes how disturbances
propagate, the new boundary between permitted and forbidden signal directions
and the possibility of controllable faster-than-light signaling would have
to be derived explicitly.

An algebraic relationship between the real and imaginary parts is shown by
the mixed terms in \(\detf{\pv{Z}}\), but a weak physical interaction is not
established. Such an interaction would require an equation of evolution and
a source law. Only then could constants to be measured experimentally be
identified. TODO questions could then be asked about a real--imaginary mixing
strength, a mass or length scale that reduces the imaginary part, or a
relative propagation or decay scale. None of these constants is predicted by
the algebra alone; each would have to be connected to a definite measured
quantity and constrained by experiment.

Allowing complex coefficients in the algebra cannot by itself be tested
experimentally because it is a mathematical extension. A proposed physical
interpretation could be rejected if measured intervals or energies became
unavoidably complex, probabilities failed to remain real and nonnegative,
time evolution failed to preserve normalization, source-free solutions grew
without bound, or predicted departures from the observed light cone and
ordinary Lorentz behavior exceeded experimental limits. The final TODO
question is whether any nontrivial complex-event evolution law can pass these
tests or whether every consistently measured result reduces to the
real-paravector theory already developed.

\subsection{FIXME: Proposed Connections to Large-Scale Observations}
No derivation is established here for energy assigned to otherwise empty
space, for the observed speeding-up of large-scale expansion, or for the
additional gravitational effects commonly attributed to unseen matter
\cite{weinberg2008}. These topics are collectively described here as
dark-sector questions. A contribution to gravity, a tendency to gather with
matter, or a measured energy is not implied by the existence of imaginary
algebra components.

A consistent local gravitational equation and a conserved physical source
must first be supplied by any such model. Ordinary local gravity must then be
recovered. The TODO question is whether quantitative predictions can be made
for the changing distance between far-separated galaxies, the bending of
light by gravity, the formation of large matter concentrations, and the
motion of matter \cite{weinberg2008}. A nearly constant contribution spread
over large distances and a contribution that gathers into dense regions are
physically different requirements and must not be inferred from the same
algebraic term without separate derivations.

\subsection{TODO: Limits of the Notation-Light Form}
The compact notation is most effective when the relevant scalar and vector
data are carried by one algebra element. In some extensions, maps with several
independent inputs, information that is valid only on overlapping local
regions, or conditions on the whole spatial region that cannot be recovered
from one local formula
are instead required. That mathematical structure is not removed by avoidance of indices;
only its recorded form is changed.

The conditions under which a paravector is sufficient, a map that preserves
addition and multiplication by ordinary numbers is required, or
several-input or whole-region data must be introduced should be stated
explicitly. Any new notation should be justified by a construction that
cannot be expressed clearly with the existing functions and operators.
The raw-text character of the framework is thereby preserved while
notational economy is prevented from being mistaken for a reduction in
physical content.
